\documentclass{article}\usepackage{amsmath}\usepackage{graphicx}\usepackage{geometry}\geometry{a4paper, margin=1in}\title{SelAction Program: Technical Description of \texttt{selovlp.f90}}\author{Generated by Gemini}\date{\today}\begin{document}\maketitle\begin{abstract}This document provides a detailed technical description of the \texttt{selovlp.f90} module from the SelAction program. This module implements selection for populations with overlapping generations.\end{abstract}\tableofcontents\section{Overlapping Generations Model}The \texttt{ovlp} subroutine in the \texttt{selovlp.f90} module handles selection in age-structured populations. This is a more complex scenario than discrete generations, as animals of different ages contribute to the next generation.\subsection{Genetic Gain per Year}The primary output is the annual genetic gain, which is calculated from the genetic gain per generation and the generation interval ($L$):\begin{equation}\Delta G_{year} = \frac{\Delta G_{gen}}{L}\end{equation}The generation interval is the weighted average of the age of the parents at the birth of their offspring. It is calculated separately for sires ($L_s$) and dams ($L_d$):\begin{equation}L = \frac{L_s + L_d}{2}\end{equation}\subsection{Gene Flow Matrix}The program uses a gene flow matrix approach to model the contribution of different age classes to the gene pool of the next generation. The response to selection is calculated for each age class and then weighted by the proportion of selected animals from that class to get the overall genetic gain.\subsection{Iterative Solution}Similar to the discrete generation model, the program iteratively solves for the equilibrium genetic parameters. The process is more complex due to the age structure. The genetic and phenotypic (co)variance matrices are updated based on the selection across all age classes. The process is repeated for 25 iterations to reach a stable solution.\end{document}